% English translation of checked_pages/esq_119_026.tex; source: cote 119, site 26, PDF page 27, typescript page 24.
% AI-translated by GPT-6 Astra (OpenAI, ChatGPT Pro mode).
\scanpage{119/26}
\mspage{24}

for the icosahedron, characteristic 2 for the octahedron), one reads
off the geometric regular polyhedron over the finite field concerned
($\mathbb{F}_4$ and $\mathbb{F}_5$ for the icosahedron, $\mathbb{F}_2$ for the octahedron) a particularly
elegant (and unexpected) description of the combinatorics
of the polyhedron. I even seemed to glimpse here a principle
of great generality, which I believed I found again, in particular, in
a later reflection on the combinatorics of the system of 27
lines on a cubic surface and its relations with the root
system $E_7$. Whether such a principle really exists and one
even succeeds in bringing it out\frapparatus{V-FR-026-001} of its cloak of mist, or whether it
recedes as one pursues it and ends by
vanishing like a Fata Morgana,\tnote{TN-026-01}{``Fata Morgana'' is retained as the author's image of an elusive mirage.}
for my part I \unsure{find}\tnote{TN-026-02}{The French alternatives remain ``trcuve|trouve''; both are represented by ``find'', with the doubt linked to the scan.}
in it an uncommon motivating force and fascination,
like that of a dream, perhaps. There is no doubt that following such a call from the
unformulated, from the formless seeking form, from an elusive
glimpse that seems to take pleasure both in evading one and in
revealing itself can only lead far, while no one could
predict where ...

Yet, occupied with other interests and tasks, I have
not so far followed this call, nor met anyone else
who wished to hear it, still less to follow it. Apart from
a few digressions toward other kinds of geometric-combinatorial
structures, my work here too has been confined to an initial
roughing-out and housekeeping, on which it is unnecessary
to dwell further here\hand{$^{(5)}$}. The only point that perhaps still deserves
mention is the existence and uniqueness of the hyperquadric circumscribed
about a given regular n-polyhedron, whose equation can be made
explicit by simple formulas in terms of the fundamental parameters
$\alpha_i$ (*). The case that interests me most is the one where $n = 2$, and the

\rule{5cm}{0.4pt}

(*) An analogous result holds for pseudo-polyhedra\unsure{.|:} It would seem
that the ``exceptional characteristics'' discussed above,
for the specializations of a given polyhedron, are those for
which the circumscribed \var{hyperquadratic}{hyperquadric}{slip corrected}{V-FR-026-002} is either degenerate or
tangent to the hyperplane at infinity.

\margin{26}

